\documentclass[12pt,a4paper]{article}
\usepackage[margin=1in]{geometry}
\usepackage[T1]{fontenc}
\usepackage[utf8]{inputenc}
\usepackage{amsmath,amssymb,amsthm}
\usepackage{microtype}

\title{A Closure--Ladder for Number Systems:\\
From $\mathbb{N}$ to $\mathbb{H}$ and Beyond}
\author{Thomas Bradley}
\date{\today}

% Theorem styles
\newtheorem{theorem}{Theorem}
\newtheorem{proposition}[theorem]{Proposition}
\newtheorem{corollary}[theorem]{Corollary}
\theoremstyle{definition}
\newtheorem{definition}[theorem]{Definition}
\theoremstyle{remark}
\newtheorem*{remark}{Remark}

\begin{document}
\maketitle

\begin{abstract}
We formalize a ``closure--ladder'' narrative: each extension of the standard number systems arises to restore closure under an operation that previously failed. We give succinct statements (with short proofs) witnessing: (i) $\mathbb{Z}$ closes subtraction absent in $\mathbb{N}$; (ii) $\mathbb{Q}$ closes division (by nonzero) absent in $\mathbb{Z}$; (iii) $\mathbb{R}$ closes Cauchy limits (completeness) absent in $\mathbb{Q}$; (iv) $\mathbb{C}$ closes roots of negative reals absent in $\mathbb{R}$; (v) $\mathbb{H}$ is \emph{necessary} to realize non-commutative multiplication (the commutator), which is trivial in $\mathbb{C}$. We also include an analytic insufficiency: $\mathbb{C}$ is not closed under the inverse of second-order tetration (the ``super-root''), since $z^z = 0$ has no complex solution.
\end{abstract}

\section{Axioms and Notation}
We regard $\mathbb{N}\subset\mathbb{Z}\subset\mathbb{Q}\subset\mathbb{R}\subset\mathbb{C}$ in the standard way and $\mathbb{H}$ as the quaternion algebra with basis $1,i,j,k$ and $i^2=j^2=k^2=ijk=-1$.
For $x,y$ in a ring, the \emph{commutator} is $[x,y]=xy-yx$.

\section{Arithmetic Ladder}
\begin{proposition}[Subtraction closure forces $\mathbb{Z}$]
\label{prop:Zclosesub}
$\mathbb{N}$ is not closed under subtraction: e.g.\ $1-2\notin\mathbb{N}$. The smallest extension containing $\mathbb{N}$ that is closed under subtraction is $\mathbb{Z}$.
\end{proposition}

\begin{proof}
$1-2=-1\notin\mathbb{N}$. Adjoin additive inverses for all $n\in\mathbb{N}$ and close under $+$ to obtain $\mathbb{Z}$; minimality is by construction.
\end{proof}

\begin{proposition}[Division closure forces $\mathbb{Q}$]
\label{prop:Qclosediv}
$\mathbb{Z}$ is not closed under division by nonzero: $1/2\notin\mathbb{Z}$. The smallest field containing $\mathbb{Z}$ (hence closed under $/$ by nonzero) is $\mathbb{Q}$.
\end{proposition}

\begin{proof}
Witness $1/2\notin\mathbb{Z}$. The field of fractions of the integral domain $\mathbb{Z}$ is $\mathbb{Q}$; minimality is standard.
\end{proof}

\section{Analytic Ladder}
\begin{theorem}[Completeness forces $\mathbb{R}$]
\label{thm:Rcomplete}
$\mathbb{Q}$ is not complete: there exist Cauchy sequences in $\mathbb{Q}$ that do not converge in $\mathbb{Q}$ (e.g.\ rationals approximating $\sqrt{2}$). The completion of $\mathbb{Q}$ under the metric topology is $\mathbb{R}$, where every Cauchy sequence converges.
\end{theorem}

\begin{proof}[Proof sketch]
The sequence of rational truncations of $\sqrt{2}$ is Cauchy but has no rational limit. The completion construction (via Cauchy sequences modulo null sequences or Dedekind cuts) yields $\mathbb{R}$ with the stated property.
\end{proof}

\begin{proposition}[Roots of negatives force $\mathbb{C}$]
\label{prop:Crootneg}
$\mathbb{R}$ is not closed under square roots of negative numbers: $\sqrt{-1}\notin\mathbb{R}$. Adjoining a root $i$ with $i^2=-1$ yields $\mathbb{C}=\mathbb{R}[i]$, which is closed under taking square roots of negative reals.
\end{proposition}

\begin{proof}
By order properties of $\mathbb{R}$, no real squares to $-1$. In $\mathbb{C}$, $\sqrt{-a}=i\sqrt{a}$ for $a>0$.
\end{proof}

\section{Structural Ladder: Necessity of $\mathbb{H}$}
\begin{theorem}[Commutator is trivial in $\mathbb{C}$, nontrivial in $\mathbb{H}$]
\label{thm:Hnecessary}
For all $z_1,z_2\in\mathbb{C}$, $[z_1,z_2]=0$ (commutativity). In $\mathbb{H}$, $[i,j]=ij-ji=k-(-k)=2k\neq 0$. Hence $\mathbb{H}$ is necessary to realize non-commutative multiplication.
\end{theorem}

\begin{proof}
Complex multiplication is commutative, so $z_1z_2=z_2z_1$. In $\mathbb{H}$ the standard relations give $ij=k$ and $ji=-k$.
\end{proof}

\begin{remark}[Interpretation]
The operation $[x,y]$ is a bona fide algebraic operator. Over $\mathbb{C}$ it collapses identically to $0$; over $\mathbb{H}$ it attains nonzero values. This is a precise sense in which $\mathbb{H}$ ``fills a structural hole'' left by $\mathbb{C}$.
\end{remark}

\section{Analytic Insufficiency: Super-root over $\mathbb{C}$}
\begin{definition}[Principal complex power]
For $z\in\mathbb{C}\setminus(-\infty,0]$, define $z^z:=\exp(z\log z)$ where $\log$ is the principal branch.
\end{definition}

\begin{theorem}[$\mathbb{C}$ is not closed under the inverse of $z\mapsto z^z$ at $w=0$]
\label{thm:superroothole}
The equation $z^z=0$ has no solution $z\in\mathbb{C}$ (on any branch), hence $w=0$ is not in the image of $z\mapsto z^z$ and the inverse operation (the second-order super-root) is not everywhere defined on $\mathbb{C}$.
\end{theorem}

\begin{proof}
For any choice of logarithm branch, $z^z=\exp\!\big(z\log z\big)$. The complex exponential never attains $0$, so $\exp(\cdot)=0$ has no finite solution. At $z=0$, one has $\lim_{z\to 0} z\log z=0$ and thus $\lim_{z\to 0} z^z=1$, not $0$. Therefore $z^z\neq 0$ for all $z\in\mathbb{C}$.
\end{proof}

\section{Synthesis: A Closure--Ladder Template}
\begin{definition}[Closure ladder step]
An extension $E\supset F$ is a \emph{closure step} for an operator $\mathcal{O}$ if $\mathcal{O}$ is not closed over $F$ (some inputs in $F$ have outputs not in $F$ or are undefined) but is closed or nondegenerate over $E$.
\end{definition}

\begin{theorem}[Exhibit of steps]
\label{thm:ladder}
The following are closure steps in the above sense:
\[
\mathbb{N}\xrightarrow{\ \text{subtraction}\ }\mathbb{Z},\quad
\mathbb{Z}\xrightarrow{\ \text{division}\ }\mathbb{Q},\quad
\mathbb{Q}\xrightarrow{\ \text{Cauchy limits}\ }\mathbb{R},\quad
\mathbb{R}\xrightarrow{\ \text{roots of negatives}\ }\mathbb{C},\quad
\mathbb{C}\xrightarrow{\ \text{noncommutative product}\ }\mathbb{H}.
\]
Moreover, the map $z\mapsto z^z$ exhibits an analytic insufficiency over $\mathbb{C}$ (Theorem~\ref{thm:superroothole}).
\end{theorem}

\begin{remark}[About the ``Quaternion Bridge'']
Finding $w\in\mathbb{C}$ such that no $z\in\mathbb{C}$ solves $z^z=w$ but some $q\in\mathbb{H}$ solves $q^q=w$ would show a \emph{tetration-driven} step from $\mathbb{C}$ to $\mathbb{H}$. Simple candidates (e.g.\ $q=j$ yield $q^q=e^{-\pi/2}$) fail because $i^i=e^{-\pi/2}$ already lies in the complex image. This remains an open research direction here.
\end{remark}

\end{document}
